\chapter{Validazione}
\label{cap:validazione}

\begin{inbreve}
La validazione stima empiricamente l'errore di generalizzazione, per due scopi distinti: scegliere il modello
(\emph{model selection}) e stimarne il rischio (\emph{model assessment}). Il capitolo presenta hold-out e K-fold
cross-validation e la ricerca degli iperparametri. Segue gli schemi formali che combinano selezione e
valutazione. Chiude con gli errori più frequenti e le misure d'errore per classificazione e regressione. Gli
aspetti specifici delle reti neurali (epoche, early stopping, inizializzazioni) sono nel
\rif{sec:validazione-reti}.
\end{inbreve}

\section{Perché validare}

La prestazione di un sistema di ML è la sua \textbf{accuratezza predittiva} su dati nuovi. Le prestazioni sui
dati di training danno una valutazione \textbf{troppo ottimistica}. All'inizio dell'addestramento c'è troppo
bias (underfitting, errore di training alto); alla fine c'è troppa varianza (errore di training basso ma
overfitting). Inoltre, a parità di complessità, l'errore di test varia molto da un training set all'altro
(\rif{sec:bv-lambda}). Un iperparametro $\theta$, implicito o esplicito, può regolare la complessità del modello.
Si vuole allora il valore di $\theta$ che minimizza l'errore di test. Servono quindi metodi per \textbf{stimare
l'errore atteso} di un modello o di una classe di modelli. Oltre ai metodi analitici (\rif{sec:rischio}) ci sono
i metodi \textbf{empirici}, basati sul ricampionamento dei dati: \textbf{cross-validation} e \textbf{bootstrap}.

\section{Model selection e model assessment}

\begin{definizione}[I due obiettivi della validazione]
Dopo aver addestrato i modelli sul training set:
\begin{itemize}
  \item \termine{model selection}: stimare le prestazioni (errore di generalizzazione) di diversi modelli per
        \textbf{scegliere il migliore}. Include la ricerca dei migliori \textbf{iperparametri} (grado del
        polinomio, numero di unità di una rete, $\lambda$, $\eta$, \dots). \emph{Restituisce un modello.}
  \item \termine{model assessment}: scelto il modello finale, \textbf{stimarne l'errore di predizione} (rischio)
        su dati di test nuovi, come misura della sua qualità. \emph{Restituisce una stima.}
\end{itemize}
\end{definizione}

\begin{attenzione}[Regola d'oro]
\uline{Tenere separati gli obiettivi e usare insiemi di dati separati in ogni fase.} Se il test set viene usato
per scegliere il modello, la stima finale non è più affidabile (è ottimistica).
\end{attenzione}

\subsection{Hold-out}

Nel mondo ideale si avrebbero un grande training set, un grande validation set e un test set esterno molto
grande. Se i dati sono sufficienti, si usa la tecnica di base, l'\termine{hold-out}. Divide il dataset in tre
insiemi \textbf{disgiunti} (ad esempio 50\% TR, 25\% VL, 25\% TS):
\begin{itemize}
  \item \textbf{training set (TR)}: per adattare il modello, cioè minimizzare il $\Remp$ della SLT
        [\emph{training}];
  \item \textbf{validation set (VL)}, o \emph{selection set}: per scegliere il modello migliore tra modelli e
        configurazioni di iperparametri diversi [\emph{model selection}];
  \item \textbf{test set (TS)}: per stimare l'errore di generalizzazione del modello finale, cioè stimare $R$
        [\emph{model assessment}]. Non si usa mai per scegliere o regolare $h$.
\end{itemize}
TR e VL insieme formano il \termine{development} o \termine{design set}, usato per costruire il modello finale
(\figrif{fig:tr-vl-ts}).

\begin{figure}[ht]
\centering
\begin{tikzpicture}
\draw[fill=grigiochiaro] (0,0) rectangle (4.6,0.7) node[midway, font=\small] {TR};
\draw (4.6,0) rectangle (7.0,0.7) node[midway, font=\small] {VL};
\draw[fill=inchiostro] (7.0,0) rectangle (9.4,0.7) node[midway, font=\small, text=white] {TS};
\draw[decorate, decoration={brace, amplitude=4pt}] (0,0.8) -- (7.0,0.8) node[midway, above=3pt, font=\footnotesize] {design set};
\node[blocco, minimum width=2.6cm] (train) at (2.3,-1.3) {training\\del modello};
\node[blocco, minimum width=2.6cm] (sel) at (5.8,-1.3) {selezione\\del modello};
\node[blocco, minimum width=2.6cm] (ass) at (9.3,-1.3) {valutazione del\\modello finale};
\node[blocco, minimum width=2.6cm, dashed] (uso) at (12.6,-1.3) {uso: predizioni\\su dati nuovi};
\draw[->] (2.3,0) -- (train); \draw[->] (5.8,0) -- (sel); \draw[->] (8.2,0) -- (ass.north west);
\draw[->] (train) -- (sel); \draw[->] (sel) -- (ass); \draw[->] (ass) -- (uso);
\end{tikzpicture}
\caption{Ruolo di TR, VL e TS: sviluppo del modello (training e selezione), stima del rischio sul test set, uso
del modello rilasciato.}
\label{fig:tr-vl-ts}
\end{figure}

\begin{attenzione}[Due precisazioni]
\begin{enumerate}
  \item La stima fatta sul VL serve \textbf{solo} alla model selection: non è una buona stima per l'assessment.
  \item I risultati sul TS \textbf{non} si possono usare per la model selection. Se lo si fa, quello non è più
        un test set: è un validation set.
\end{enumerate}
\end{attenzione}

Si pensi a un ciclo di progettazione ripetuto: provo, guardo il test, cambio, riprovo. Usare il test set in
questo modo è \textbf{model selection}, non una valutazione. Sugli esempi futuri non si potrebbe fare lo stesso, e
l'errore sul test così usato è una \textbf{stima troppo ottimistica}. È come un esercizio d'esame di cui si sono
viste le soluzioni. Da qui il concetto di \termine{blind test set} (test cieco), usato nelle competizioni.

\section{Un controesempio istruttivo}
\label{sec:controesempio}

\begin{esempio}[Il target casuale]
Si hanno 20--30 esempi con \textbf{1000 variabili di input} a valori casuali e un target \textbf{casuale}
$0/1$. Si seleziona il modello che usa \textbf{una sola} variabile di input. È quella che per puro caso
``indovina'' il target al 99\% (o 100\%) su qualsiasi successiva suddivisione in TR, VL e TS. Risultato perfetto?
No: la vera accuratezza è il \textbf{50\%}, perché il target è casuale e nessun modello può fare meglio di una
moneta. Su un test set davvero nuovo la stessa variabile darebbe accuratezze del tutto casuali (ad esempio 100\%,
33\%, 66\% su tre test set diversi).
\end{esempio}

Con 1000 variabili casuali e soli 20 esempi è quasi certo che una variabile coincida per caso con il target sulla
maggior parte degli esempi. La scelta è fatta guardando \textbf{tutti} i dati. Per questo quella variabile
``funziona'' su ogni partizione successiva. Gli errori sono due:
\begin{enumerate}
  \item[(1)] la stima dell'errore su TR e VL \textbf{non è} una buona stima del rischio;
  \item[(2)] usare \textbf{l'intero dataset} per la selezione delle feature o del modello \textbf{pregiudica la
        stima} (\ingl{subset selection bias}): il test set è stato usato \emph{implicitamente} all'inizio. Il
        test set va separato \textbf{prima} di qualsiasi model selection, \textbf{feature selection compresa}.
\end{enumerate}

\begin{attenzione}
Qui si discute la \textbf{correttezza della stima}, non la possibilità di risolvere il task. La K-fold CV e le
altre tecniche \textbf{non} risolvono il problema se la selezione è fatta prima della suddivisione.
\end{attenzione}

\section{K-fold cross-validation}
\label{sec:kfold}

L'hold-out richiede dati sufficienti per tre partizioni significative. Quanti ne servano dipende dalla
complessità del modello e dal rapporto segnale/rumore. Il risultato dipende inoltre dalla \textbf{particolare
partizione} scelta. Con pochi dati l'hold-out sfrutta male l'informazione disponibile.

\begin{definizione}[K-fold cross-validation]
Si divide il dataset $D$ in $K$ sottoinsiemi mutuamente esclusivi $D_1,\dots,D_K$, i \termine{fold}. Per ogni
$i$ si addestra su $D\setminus D_i$ e si valuta su $D_i$. Infine si combinano i $K$ risultati.
\end{definizione}

\begin{figure}[ht]
\centering
\begin{tikzpicture}[x=1.15cm, y=0.62cm]
\foreach \r in {1,...,4} {
  \foreach \c in {1,...,4} {
    \ifnum\r=\c \fill[inchiostro] (\c,-\r) rectangle ++(1,0.8); \node[white, font=\footnotesize] at (\c+0.5,-\r+0.4) {$D_\c$};
    \else \draw (\c,-\r) rectangle ++(1,0.8); \node[font=\footnotesize] at (\c+0.5,-\r+0.4) {$D_\c$}; \fi
  }
  \node[font=\footnotesize, anchor=west] at (5.2,-\r+0.4) {iterazione \r};
}
\end{tikzpicture}
\caption{4-fold cross-validation: a turno ogni fold (nero) fa da insieme di validazione (o di test) e gli altri
da training.}
\label{fig:kfold}
\end{figure}

Così tutti i dati vengono usati sia per l'addestramento sia per la valutazione. Non si ottiene un \textbf{unico
modello}: ce n'è uno per fold. Si ottiene però anche una \textbf{varianza} dei risultati sui fold, e il risultato
si riporta come \textbf{media $\pm$ deviazione standard}. Si può usare sia per il validation set sia per il test
set. Restano tre questioni aperte. La prima è il numero di fold: 3, 5, 10, fino al \termine{leave-one-out}, con
$K=l$. La seconda è il costo computazionale. La terza è la combinazione con altri schemi
(\rif{sec:schemi-combinati}).

Il termine ``cross-validation'' ha due significati. Indica un \textbf{approccio} alla model selection, basato sulla
stima diretta dell'errore tramite ricampionamento (anziché analitica). Indica anche le \textbf{implementazioni}
che stabiliscono come dividere i dati, come la K-fold CV. Un'altra tecnica di ricampionamento è il
\termine{bootstrap}: un campionamento casuale \textbf{con reinserimento}, ripetuto per ottenere sottoinsiemi
diversi di validazione o di test.

\section{Ricerca degli iperparametri}
\label{sec:grid-search}

Gli \termine{iperparametri} sono i parametri non appresi direttamente dall'algoritmo di apprendimento. Si fissano
cercando nello spazio degli iperparametri, tramite model selection su un validation set.

\subsection{Grid search}

La \termine{grid search} esaustiva genera tutti i candidati di una griglia di valori. Ad esempio con numero di
unità e $\lambda$:

\begin{center}\small
\begin{tabular}{l|ccc}
Unità $\backslash$ $\lambda$ & $0{,}1$ & $0{,}01$ & $0{,}001$\\
\hline
1 unità & Res1 & Res4 & Res7\\
10 unità & Res2 & Res5 & Res8\\
100 unità & \textbf{Res3} & Res6 & Res9\\
\end{tabular}
\end{center}
Se il risultato migliore sul validation set è Res3, vince la configurazione (100 unità, $\lambda=0{,}1$). La
ricerca va \textbf{automatizzata} ed è facile da \textbf{parallelizzare}, perché le prove sono indipendenti. Il
costo è il prodotto cartesiano dei valori, $(\#\text{valori})^{\#\text{iperparametri}}$. Con 3 valori per 2
iperparametri sono $3^2=9$ prove; con 5--6 iperparametri da 10 valori il costo diventa proibitivo. Per questo:
\begin{itemize}
  \item si possono fissare in una fase preliminare gli iperparametri che risultano poco rilevanti (con cautela,
        \rif{sec:errori-validazione});
  \item si usano \textbf{grid search annidate}. Prima si fa una ricerca \textbf{grossolana} su tutte le
        combinazioni, ad esempio con valori a crescita esponenziale, per trovare le regioni buone. Poi si fanno
        ricerche \textbf{più fini} su intervalli sempre più piccoli. Si considerano sempre \textbf{tutti} gli
        iperparametri significativi insieme.
\end{itemize}

\subsection{Alternative}

La grid search è molto utile per \textbf{osservare il comportamento} del modello al variare degli iperparametri,
ma scala male. Le alternative sono la \termine{random search} e la ricerca \textbf{automatica}: approcci
bayesiani, evolutivi, metodi come Hyperband, verso l'\ingl{AutoML}. Le librerie ne includono molte, ma non vanno
usate acriticamente.

\begin{idea}[Perché la random search funziona]
Si supponga che solo uno di due iperparametri conti davvero. Una griglia $10\times10$ ne prova appena \textbf{10
valori distinti}: le altre 90 prove ripetono gli stessi valori, variando l'iperparametro inutile. 100 prove
casuali ne provano invece \textbf{100} (\figrif{fig:grid-random}). La ricerca casuale permette anche di fissare il
\textbf{budget} di prove indipendentemente dal numero di iperparametri.
\end{idea}

\begin{figure}[ht]
\centering
\begin{tikzpicture}
\begin{groupplot}[group style={group size=2 by 1, horizontal sep=1.0cm}, scatola, width=5.4cm, height=5.4cm,
  xmin=0, xmax=1, ymin=0, ymax=1, xtick=\empty, ytick=\empty, xlabel={iperparametro importante},
  ylabel={iperparametro irrilevante}]
\nextgroupplot[title={\small grid search}]
\addplot[only marks, mark=*, mark size=1.1pt, inchiostro] table {dati/grid_griglia.dat};
\addplot[only marks, mark=|, mark size=4pt, verde, thick] table[y expr=0] {dati/grid_griglia.dat};
\nextgroupplot[title={\small random search}]
\addplot[only marks, mark=*, mark size=1.1pt, inchiostro] table {dati/grid_random.dat};
\addplot[only marks, mark=|, mark size=4pt, verde, thick] table[y expr=0] {dati/grid_random.dat};
\end{groupplot}
\end{tikzpicture}
\caption{100 prove con grid search e con random search: in verde, sull'asse orizzontale, i valori distinti
provati per l'iperparametro importante (10 contro 100).}
\label{fig:grid-random}
\end{figure}

\section{Schemi formali}
\label{sec:schemi}

Una notazione precisa evita ambiguità e mostra dove verrebbero commessi gli errori del controesempio. Gli schemi
valgono per la selezione tra configurazioni di iperparametri $\theta$, ma anche tra \textbf{modelli diversi}:
$\theta$ può indicare il tipo di modello.

\begin{definizione}[Notazione]
\begin{itemize}
  \item \textbf{Dati}: $z_1,\dots,z_l$ campione i.i.d.\ di una distribuzione sconosciuta $p(z)$;
        $D_r = \{z_1,\dots,z_r\}$ è un insieme di $r$ elementi.
  \item \textbf{Task}: ogni dato è una coppia $z = (\x,d)$. Per la classificazione $z\in\R^n\times\{-1,+1\}$, per
        la regressione $z\in\R^n\times\R$.
  \item \textbf{Loss}: classificazione $L((\x,d),h) = 0$ se $h(\x)=d$, $1$ altrimenti; regressione
        $L((\x,d),h) = (h(\x)-d)^2$.
  \item \textbf{Rischio atteso} e ipotesi ottima:
        $R(h) = \E_z[L(z,h)] = \int L(z,h)\,p(z)\,dz$, \quad $h^* = \argmin_{h\in H}R(h)$.
  \item \textbf{Rischio empirico} su un insieme finito: $\Remp(h,D_r) = \frac1r\sum_{i=1}^{r}L(z_i,h)$.
\end{itemize}
\end{definizione}

\begin{attenzione}
Qui $\Remp$ è inteso in senso generale, come \textbf{surrogato di $R$} su un insieme finito. Si usa per il
training, la selezione o la valutazione, a seconda della partizione. Il $\Remp$ della SLT (l'errore di training)
è, in questa notazione, $\Remp(h,D_{TR})$.
\end{attenzione}

Prima si identifica l'obiettivo. Può essere il \textbf{miglior modello} ottenibile dal dataset (model
selection). Può essere la stima delle \textbf{prestazioni attese} di un modello ottenuto minimizzando il rischio
empirico (model assessment). Oppure possono servire \textbf{entrambi}.

\subsection{Model selection}

\begin{algoritmo}[Model selection con hold-out]
\begin{enumerate}
  \item Scegliere uno spazio delle ipotesi con iperparametro $\theta$; dividere $D_l$ in
        $D_{TR} = \{z_1,\dots,z_{tr}\}$ e $D_{VL} = \{z_{tr+1},\dots,z_{tr+vl}\}$, con $tr+vl=l$.
  \item Per ogni valore $\theta_m$ [grid search]:
        \begin{itemize}
          \item $h^*_{\theta_m}(D_{TR}) = \argmin_{h\in H_{\theta_m}}\Remp(h,D_{TR})$ [training];
          \item stimare $R(h^*_{\theta_m})$ con
                $\Remp(h^*_{\theta_m},D_{VL}) = \frac{1}{vl}\sum_{z_i\in D_{VL}}L\big(z_i,h^*_{\theta_m}(D_{TR})\big)$.
        \end{itemize}
  \item Scegliere $\theta^*_m = \argmin_{\theta_m}R(h^*_{\theta_m})$ [miglior risultato sul VL].
  \item Restituire $h^*(D_l) = \argmin_{h\in H_{\theta^*_m}}\Remp(h,D_l)$ [riaddestramento su tutti gli $l$ dati].
\end{enumerate}
\end{algoritmo}

La stima di $R$ sul VL serve solo alla scelta e \textbf{non viene restituita} come stima del rischio: farlo
sarebbe l'errore (1) del controesempio. Il motivo è che quella stima è ottimistica. Ogni errore di validazione è
affetto da rumore, e si è scelto proprio il più basso tra molti: il minimo tende a essere stato ``fortunato''.

\begin{algoritmo}[Model selection con K-fold CV]
\begin{enumerate}
  \item Scegliere uno spazio delle ipotesi con iperparametro $\theta$; dividere $D_l$ in $K$ parti distinte e
        uguali $D_1,\dots,D_K$; sia $\bar D_k$ il complementare di $D_k$.
  \item Per ogni valore $\theta_m$ [grid search]:
        \begin{itemize}
          \item per ogni parte $D_k$ [ciclo sui fold]: $h^*_{\theta_m}(\bar D_k) = \argmin_{h\in H_{\theta_m}}\Remp(h,\bar D_k)$
                [training], e stima di $R(h^*_{\theta_m}(\bar D_k))$ con
                $\Remp(h^*_{\theta_m}(\bar D_k),D_k) = \frac{1}{\abs{D_k}}\sum_{z_i\in D_k}L\big(z_i,h^*_{\theta_m}(\bar D_k)\big)$;
          \item stima di $R(h^*_{\theta_m}(D_l))$ con la media $\frac1K\sum_kR(h^*_{\theta_m}(\bar D_k))$.
        \end{itemize}
  \item Scegliere $\theta^*_m = \argmin_{\theta_m}R(h^*_{\theta_m}(D_l))$ [migliore media di validazione].
  \item Restituire $h^*(D_l) = \argmin_{h\in H_{\theta^*_m}}\Remp(h,D_l)$ [riaddestramento].
\end{enumerate}
\end{algoritmo}

\begin{attenzione}[Ogni iterazione riparte da zero]
Addestramento e stima sono locali a ciascuna suddivisione. Si supponga di usare il modello addestrato nello split
1 per valutare lo split 2. Si userebbe come validazione dello split 2 una parte ($D_2$) già usata per addestrare
nello split 1.
\end{attenzione}

\subsection{Model assessment}

\begin{algoritmo}[Model assessment con hold-out]
\begin{enumerate}
  \item Dividere $D_l$ in $D_{TR}$ e $D_{TS}$, con $tr+ts=l$. \uline{Il test set va separato adesso}, non dopo:
        altrimenti si commette l'errore (2) del controesempio.
  \item $h^*(D_{TR}) = \argmin_{h\in H}\Remp(h,D_{TR})$ [training; può includere una model selection interna su
        $D_{TR}$].
  \item Stimare $R(h^*(D_{TR}))$ con $\Remp(h^*(D_{TR}),D_{TS}) = \frac{1}{ts}\sum_{z_i\in D_{TS}}L\big(z_i,h^*(D_{TR})\big)$
        [errore di test].
\end{enumerate}
\end{algoritmo}

\begin{algoritmo}[Model assessment con K-fold CV]
\begin{enumerate}
  \item Dividere $D_l$ in $K$ parti distinte e uguali.
  \item Per ogni parte $D_k$: $h^*(\bar D_k) = \argmin_{h\in H}\Remp(h,\bar D_k)$ [training, eventualmente con
        model selection], e stima di $R(h^*(\bar D_k))$ con $\Remp(h^*(\bar D_k),D_k)$ [errore di test della
        iterazione].
  \item Stimare $R(h^*(D_l))$ con la media $\frac1K\sum_kR(h^*(\bar D_k))$.
\end{enumerate}
\end{algoritmo}

Con $K=l$ si ha la \textbf{leave-one-out CV}. Non si violano i principi: i risultati di test non vengono mai usati
per addestrare o selezionare, in nessuna iterazione.

\subsection{Combinare selezione e valutazione}
\label{sec:schemi-combinati}

Quando servono sia il modello migliore sia la stima del suo rischio si combinano gli schemi.

\begin{algoritmo}[Hold-out TR--VL--TS]
\begin{enumerate}
  \item Dividere $D_l$ in $D_{TR}$, $D_{VL}$, $D_{TS}$.
  \item Per ogni $\theta_m$: addestrare $h^*_{\theta_m}(D_{TR})$ e calcolare $\Remp(h^*_{\theta_m}(D_{TR}),D_{VL})$.
  \item Scegliere $\theta^*_m$ con il miglior risultato sul VL.
  \item Riaddestrare $h^*(D_{TR}\cup D_{VL})$ con $\theta^*_m$ (eventualmente usando una parte dei dati come VL
        per l'early stopping).
  \item Stimare $R$ con $\frac{1}{ts}\sum_{z_i\in D_{TS}}L\big(z_i,h^*(D_{TR}\cup D_{VL})\big)$.
\end{enumerate}
\end{algoritmo}

\begin{algoritmo}[K-fold CV per la selezione e hold-out per il test]
\begin{enumerate}
  \item Dividere $D_l$ in $D_{TR}$ (il design set) e $D_{TS}$.
  \item Per ogni $\theta_m$: stimare $R(h^*_{\theta_m}(D_{TR}))$ con una K-fold CV su $D_{TR}$ (errore medio di
        validazione sui fold).
  \item Scegliere $\theta^*_m$ con la stima migliore; riaddestrare $h^*(D_{TR})$.
  \item Stimare $R$ con l'errore su $D_{TS}$.
\end{enumerate}
\end{algoritmo}

Ecco un esempio completo di questo schema. Si divide il dataset in TR e TS. Per ogni cella di una grid search
(ad esempio $\lambda=0{,}1$ e 20 unità) si esegue un'intera K-fold CV dentro il TR. Si sceglie la configurazione
con il miglior errore medio di validazione. Si riaddestra sull'intero TR e si valuta sul TS esterno.

\subsubsection{K-fold esterna per il test, hold-out interno per la selezione}
In ogni iterazione della K-fold esterna un fold fa da test. Il resto viene diviso a caso (con shuffling) in TR e
VL per la model selection (\figrif{fig:kfold-holdout}). Questo processo \textbf{non fornisce un modello finale}:
ogni riga può selezionare un modello diverso. Si ottiene solo una stima del rischio della \textbf{classe} di
modelli considerata.

\begin{figure}[ht]
\centering
\begin{tikzpicture}[x=1.15cm, y=0.62cm]
\foreach \r in {1,...,4} {
  \foreach \c in {1,...,4} {
    \ifnum\r=\c \fill[inchiostro] (\c,-\r) rectangle ++(1,0.8); \node[white, font=\scriptsize] at (\c+0.5,-\r+0.4) {TS};
    \else \draw (\c,-\r) rectangle ++(1,0.8);
      \pgfmathtruncatemacro{\ultimo}{\r==4 ? 3 : 4}
      \ifnum\c=\ultimo \fill[grigiochiaro] (\c,-\r) rectangle ++(1,0.8); \draw (\c,-\r) rectangle ++(1,0.8); \node[font=\scriptsize] at (\c+0.5,-\r+0.4) {VL};
      \else \node[font=\scriptsize] at (\c+0.5,-\r+0.4) {TR}; \fi
    \fi
  }
}
\end{tikzpicture}
\caption{K-fold CV esterna per il test, con divisione interna TR/VL per la model selection (il riaddestramento non
è mostrato).}
\label{fig:kfold-holdout}
\end{figure}

\subsubsection{Double (nested) K-fold CV}

\begin{algoritmo}[Double K-fold CV]
\begin{enumerate}
  \item Dividere $D_l$ in $K$ parti distinte e uguali.
  \item Per ogni parte $D_k$ [ciclo esterno]:
        \begin{itemize}
          \item scegliere $h^*(\bar D_k)$ con una \textbf{K-fold CV interna} su $\bar D_k$ (con $K'$ eventualmente
                diverso da $K$) [ciclo interno];
          \item stimare $R(h^*(\bar D_k))$ con $\Remp(h^*(\bar D_k),D_k)$ [errore di test di ogni split esterno].
        \end{itemize}
  \item Stimare $R(h^*(D_l))$ con la media sui fold di test.
\end{enumerate}
\end{algoritmo}

\begin{figure}[ht]
\centering
\begin{tikzpicture}[x=0.95cm, y=0.6cm]
\node[font=\small, anchor=west] at (0,1.2) {resampling esterno (stima delle prestazioni)};
\foreach \c in {1,...,4} {
  \ifnum\c=2 \fill[inchiostro] (\c-1,0) rectangle ++(1,0.8); \node[white, font=\scriptsize] at (\c-0.5,0.4) {TS};
  \else \draw (\c-1,0) rectangle ++(1,0.8); \node[font=\scriptsize] at (\c-0.5,0.4) {TR}; \fi
}
\draw[->, thick] (2,-0.3) -- (2,-1.1) node[midway, right, font=\footnotesize] {sul TR esterno};
\node[font=\small, anchor=west] at (-0.1,-1.6) {resampling interno (selezione degli iperparametri)};
\foreach \r in {1,2,3} {
  \foreach \c in {1,2,3} {
    \ifnum\r=\c \fill[grigiochiaro] (\c*1.333-1.333,-2.2-\r) rectangle ++(1.333,0.8); \draw (\c*1.333-1.333,-2.2-\r) rectangle ++(1.333,0.8); \node[font=\scriptsize] at (\c*1.333-0.667,-1.8-\r) {VL};
    \else \draw (\c*1.333-1.333,-2.2-\r) rectangle ++(1.333,0.8); \node[font=\scriptsize] at (\c*1.333-0.667,-1.8-\r) {TR}; \fi
  }
}
\node[font=\footnotesize, align=left, anchor=west] at (4.5,-3.6) {miglior $\theta$ $\to$ riaddestramento\\sul TR esterno $\to$ errore sul TS};
\end{tikzpicture}
\caption{Double CV: per ogni split esterno (in alto, uno dei quattro), una CV interna sul training esterno
seleziona gli iperparametri; il ``test set'' della CV interna è in realtà un validation set.}
\label{fig:double-cv}
\end{figure}

Anche la double CV \textbf{non dà un modello finale unico}: ogni ciclo esterno può scegliere iperparametri
diversi. Dà una stima del rischio della classe di modelli, cioè dell'algoritmo con la sua procedura di selezione,
con la sua \textbf{varianza} sui fold. Se serve un modello finale si esegue una \textbf{model selection separata}
(hold-out o K-fold CV) con riaddestramento. Non si violano le regole: la stima di test è già stata fatta, e i
risultati di test non vengono mai usati per la selezione. Non è quindi vero in assoluto che ``il test viene
sempre per ultimo'': qui avviene \emph{prima} dell'ultima model selection. La double CV sfrutta meglio tutti i
dati, sia per la selezione sia per la valutazione. Il prezzo è un costo computazionale elevato.

\begin{idea}[Si può riaddestrare su tutto $D_l$?]
Sì: il riaddestramento finale su tutti i dati (TR+VL+TS) non usa i risultati del test per alcuna scelta. Il
modello finale avrà in generale prestazioni pari o migliori della stima, perché ha più dati. La stima resta
valida come stima, leggermente pessimistica, del rischio della procedura. L'importante è che il test non abbia
influenzato nessuna decisione.
\end{idea}

\section{Errori frequenti e buone pratiche}
\label{sec:errori-validazione}

\subsection{Riprogettare in base al test}
Riprogettare il modello in base al risultato sul test \textbf{non è corretto}, a meno di disporre di nuovi dati
per un nuovo test set. Ripetendo il procedimento sugli stessi dati si può ottenere qualunque risultato: il
modello si adatta a quel dataset, e \textbf{non si sta più stimando il rischio}. Non bisogna regolare il
predittore sulla base dell'errore di test, né a mano né automaticamente: sugli esempi futuri non si potrebbe
fare. La valutazione su un test interno resta utile per \textbf{verificare il proprio approccio}. Per fare un buon
modello serve una buona model selection. Per stimare le prestazioni reali serve una stima rigorosa sul test.

\subsection{La cross-validation può andare in overfitting}
Un uso intensivo della CV può portare a overfitting. Con 20--30 record e 1000 attributi si provano 1000 modelli
di regressione lineare, con un solo attributo ciascuno. Il migliore dei 1000 sembra buono sul validation set. Ma
lo sarebbe sembrato anche con un output \textbf{puramente casuale}. Un secondo esempio: si approssima un target
casuale e si usa una K-fold CV per il test. Prima però si scelgono le feature per massima correlazione con il
target, guardando \textbf{tutti} gli esempi. Si trova allora un errore del 3\% invece del vero 50\%. Rimedi:
\begin{itemize}
  \item tenere da parte un \textbf{test set aggiuntivo} prima di qualsiasi model selection, e verificare che il
        modello migliore funzioni anche su di esso;
  \item per la selezione con CV, usare ``CV + test'' o la \textbf{double CV} (rifacendo la selezione per ogni fold);
  \item per la valutazione con CV, fare la \textbf{feature selection dentro ogni fold}: si ottiene la stima
        corretta, circa il 50\% sul task casuale.
\end{itemize}

\subsection{Selezione sequenziale degli iperparametri}
\begin{attenzione}
Selezionare gli iperparametri \textbf{uno alla volta}, in sequenza, introduce un \textbf{bias legato all'ordine}.
Un esempio è scegliere il miglior $\eta$ con 10 unità e poi usarlo anche con 100 unità.
\end{attenzione}
Il procedimento corretto:
\begin{enumerate}
  \item prove iniziali per determinare gli \textbf{intervalli} giusti (ad esempio guardando la stabilità del
        training) e gli iperparametri più \textbf{influenti} (ad esempio guardando la varianza dei risultati);
  \item ricerca su una griglia con \textbf{tutte le combinazioni} degli iperparametri;
  \item eventuali griglie annidate, da grossolane a fini, sempre con \textbf{tutti gli iperparametri significativi
        insieme}, per tenere conto degli effetti incrociati.
\end{enumerate}
Bisogna però ragionare e selezionare gli aspetti rilevanti, senza affidarsi solo ad approcci automatici costosi.

\subsection{Quale cross-validation}

\begin{table}[ht]
\centering\small
\caption{Confronto tra le tecniche (con $K$ fold).}
\begin{tabularx}{0.95\linewidth}{lXX}
\toprule
 & Svantaggi & Vantaggi \\
\midrule
Hold-out (test set) & varianza: stima inaffidabile delle prestazioni future (fortuna/sfortuna) & economico \\
Leave-one-out & costoso; comportamenti anomali & non spreca dati \\
10-fold & spreca il 10\% dei dati; 10 volte più costoso dell'hold-out & spreca solo il 10\%; 10 volte più costoso invece di $l$ volte \\
3-fold & spreca più dati della 10-fold & leggermente meglio dell'hold-out \\
$K$-fold con $K=l$ & identico al leave-one-out & \\
\bottomrule
\end{tabularx}
\end{table}

La \textbf{5- o 10-fold CV} è un buon compromesso tra bias e varianza della stima, ed è raccomandata rispetto alla
leave-one-out. L'hold-out semplice è molto usato (a volte il test set è dato separatamente), ma richiede
attenzione con pochi dati.

\subsection{Esempi di model selection}
\begin{itemize}
  \item \textbf{Modelli diversi}: si sceglie la classe di modelli con il miglior errore di CV (ad esempio
        10-fold). La si riaddestra su tutti i dati, e quello è il modello da usare. Si guarda \textbf{solo
        l'errore di validazione}, non quello di training né una loro media.
  \item \textbf{Numero di unità nascoste}: l'errore di training decresce sempre con il numero di unità. Quello di
        CV ha un minimo (in pratica si provano valori come 10, 20, 30, \dots).
  \item \textbf{$k$ del K-NN}: si usa la \textbf{leave-one-out}. Per i metodi non parametrici costa quanto una
        normale predizione, perché non c'è un modello da addestrare. Non c'è garanzia che un ottimo locale di $k$
        rispetto alla LOOCV sia globale: la relazione può essere irregolare. Per esplorare molti valori si usano
        valori a crescita esponenziale ($k = 1,2,4,8,16,\dots$) o una ricerca \ingl{hill-climbing}.
\end{itemize}

\subsection{Campionamento e pochi dati}
\begin{itemize}
  \item \termine{Stratificazione}: in ogni partizione (TR e TS, o i fold) ogni classe è rappresentata circa nelle
        stesse proporzioni del dataset completo.
  \item Controllare la composizione dei fold: l'\textbf{ordine dei dati} può avere un significato (dati ordinati
        per classe o per tempo).
  \item \textbf{Hold-out o CV ripetuti} con campionamenti casuali diversi, mediando i risultati. Sono utili
        soprattutto per \textbf{confrontare modelli diversi}.
  \item Con pochissimi dati bisogna evitare, o almeno tenere in conto, alcune situazioni: classi o feature
        mancanti nel training, classi speciali non campionate in TR o TS, outlier che influenzano la media dei
        risultati, la selezione di soli casi ``facili''.
  \item Anche un \textbf{blind test set} può essere fuorviante. Succede se proviene da una distribuzione diversa,
        se è misurato con scale o tolleranze diverse, se non è pre-elaborato come i dati di training. Succede anche
        se richiede \textbf{estrapolazione} (dati fuori range).
\end{itemize}

\section{Misure d'errore}
\label{sec:misure}

\subsection{Classificazione}

Per un classificatore binario si confrontano classe reale e classe predetta nella \termine{matrice di confusione}:

\begin{center}\small
\begin{tabular}{l|cc}
Reale $\backslash$ Predetto & Positivo & Negativo\\
\hline
Positivo & TP (veri positivi) & FN (falsi negativi)\\
Negativo & FP (falsi positivi) & TN (veri negativi)\\
\end{tabular}
\end{center}

\begin{itemize}
  \item \termine{Accuratezza} $= \dfrac{TP+TN}{\text{totale}}$: frazione di esempi classificati correttamente;
  \item \termine{Sensibilità} (\ingl{True Positive rate}, \ingl{recall}) $= \dfrac{TP}{TP+FN}$: quanti positivi
        veri vengono riconosciuti;
  \item \termine{Specificità} (\ingl{True Negative rate}) $= \dfrac{TN}{FP+TN} = 1-FPR$: quanti negativi veri
        vengono riconosciuti;
  \item \termine{Precisione} $= \dfrac{TP}{TP+FP}$: quanti dei predetti positivi sono davvero positivi.
\end{itemize}
Un falso positivo equivale a un \textbf{falso allarme}.

\begin{attenzione}[L'accuratezza da sola può ingannare]
\begin{itemize}
  \item Nella classificazione binaria, un'accuratezza del 50\% è quella di un predittore \textbf{casuale}.
  \item Con \textbf{dati sbilanciati} (ad esempio 99\% di positivi) il classificatore banale che risponde sempre
        ``positivo'' ha il 99\% di accuratezza senza aver imparato nulla.
\end{itemize}
\end{attenzione}

La \termine{curva ROC} mostra il True Positive rate (sensibilità) in funzione del False Positive rate
($1-$specificità), al variare della \textbf{soglia di decisione} del classificatore (\figrif{fig:roc}). La
diagonale corrisponde al classificatore casuale. Le curve migliori salgono rapidamente verso l'angolo in alto a
sinistra e hanno un'\termine{AUC} (\ingl{Area Under the Curve}) più alta. Il classificatore ideale ha AUC $=1$,
quello casuale $0{,}5$.

\begin{figure}[ht]
\centering
\begin{tikzpicture}
\begin{axis}[stile, width=5.6cm, height=5.6cm, xmin=0, xmax=1, ymin=0, ymax=1.02, xlabel={FP rate ($1-$specificità)},
  ylabel={TP rate (sensibilità)}, xtick={0,0.5,1}, ytick={0,0.5,1}, xticklabels={$0$,$0{,}5$,$1$}, yticklabels={$0$,$0{,}5$,$1$}, domain=0:1]
\addplot[thick, grigio, densely dashed] {x} node[pos=0.75, below right, font=\footnotesize] {casuale};
\addplot[thick, blu] {x^0.45} node[pos=0.55, below right, font=\footnotesize] {};
\addplot[thick, rosso] {x^0.15};
\node[font=\footnotesize, rosso] at (axis cs:0.25,0.95) {migliore};
\node[font=\footnotesize, blu] at (axis cs:0.55,0.62) {buono};
\end{axis}
\end{tikzpicture}
\caption{Curve ROC: la diagonale è il classificatore casuale (AUC $=0{,}5$); più la curva sale verso l'angolo in
alto a sinistra, maggiore è l'area sottesa.}
\label{fig:roc}
\end{figure}

\subsection{Regressione}

Con errore $e_i = d_i - o_i$ ($d$ target, $o$ uscita del modello):
\begin{itemize}
  \item \textbf{MSE} $=\mathrm{mean}_i[e_i^2]$ e la sua radice, \textbf{RMS} (o $S$);
  \item \termine{errore assoluto medio} (MAE) $=\mathrm{mean}_i[\abs{e_i}]$;
  \item \textbf{errore assoluto massimo} $=\max_i\abs{e_i}$;
  \item \termine{coefficiente di correlazione} $r$ e \termine{coefficiente di determinazione} $r^2$ (scritto anche
        $R^2$). Ad esempio $r = \sqrt{1-S^2/S_d^2}$, con $S_d^2 = \mathrm{mean}_i[(d_i-\bar d)^2]$ varianza del
        target. Misura il grado di dipendenza lineare tra $d$ e $o$, in $[0,1]$, dove 1 è il meglio;
  \item grafici delle \textbf{uscite contro i target}; test statistici e analisi di significatività.
\end{itemize}

\begin{puntichiave}
\begin{itemize}
  \item Model selection (restituisce un modello) e model assessment (restituisce una stima): obiettivi separati,
        insiemi separati.
  \item Il risultato sul TR non stima quello su VL e TS; il risultato sul VL non stima quello sul TS.
  \item Il test set si separa prima di ogni scelta, feature selection compresa.
  \item K-fold CV: media $\pm$ deviazione standard sui fold; ogni iterazione riparte da zero.
  \item Schemi combinati: TR--VL--TS, CV + test, double CV (stima del rischio della classe, nessun modello unico).
\end{itemize}
\end{puntichiave}

\begin{domande}
\begin{enumerate}
  \item Differenza tra model selection e model assessment; perché servono insiemi separati?
  \item Descrivere il controesempio del target casuale: quali sono i due errori?
  \item Cos'è la grid search? Come se ne riduce il costo? Vantaggi della random search.
  \item Scrivere gli algoritmi di model selection e di model assessment con K-fold CV.
  \item Descrivere la double (nested) CV. Restituisce un modello?
  \item La cross-validation può andare in overfitting? Come evitarlo?
  \item Definire accuratezza, sensibilità, specificità, precisione; cos'è la curva ROC?
\end{enumerate}
\end{domande}
